% VAL-017 -- generado por herramientas/gen_casos.py; editar el texto libremente.
\section{VAL-017: P-Delta en voladizo (N1, $B_1/B_2$)}\label{sec:VAL-017}
\begin{ficha}
Fuente & [S] catálogo MAM-253 y \cite{aisc16}\\
Ticket & MAM-104\\
Engine & commit \texttt{\EngineCommit}\\
Binario & SHA-256 \texttt{\EngineBinSHA}\\
Entrada & SHA-256 \texttt{\EntradaSHA{VAL-017}}\\
Tolerancia & $|c-r|\le A+R|r|$, $A=0$, $R=0.01$ (fijada a priori)\\
Estado de datos & \directlua{mamut.estado('VAL-017')}\\
\end{ficha}

\paragraph{Modelo.}
\begin{center}\diag{voladizo}\end{center}

\paragraph{Referencia y derivación.} Derivación propia pendiente de redactar a partir del fixture del caso (provenance del Engine). % DERIV-PENDIENTE


\paragraph{Reproducción.}
\reproducir{VAL-017}

% Q2-PLACEHOLDER INICIO (VAL-017)
% Solo Q2 escribe datos/VAL-017/opensees.csv y el texto de modelo OpenSees de este bloque,
% en su propio PR. J2 no lee ni escribe valores de OpenSees.
\paragraph{Modelo de contraste (OpenSees).} \QdosPendiente
% Q2-PLACEHOLDER FIN (VAL-017)

\paragraph{Resultados.}
\begin{center}
\begin{tikzpicture}
\begin{axis}[grafico, xlabel={$P$}, ylabel={$\Delta$}]
\directlua{mamut.serie('VAL-017','referencia','x','y','ref','Referencia analítica')}
\directlua{mamut.serie('VAL-017','engine','x','y','eng','MAMUT Engine')}
\directlua{mamut.serie('VAL-017','opensees','x','y','ops','OpenSees (contraste)')}
\end{axis}
\end{tikzpicture}
\end{center}
\begin{center}\small
\begin{tabular}{rrrrrrl}\toprule
$P$ & Ref. & Engine & OpenSees & Error rel. & Tol. & Veredicto\\\midrule
\directlua{mamut.tabla('VAL-017','x','y',0,0.01)}
\bottomrule\end{tabular}
\end{center}
